\section*{Hoofdstuk \ref{ch:g6:water-underground} --- Water onder de grond}

\begin{solution}{exo:g6:water-underground:1}
\omterm{def:g6:water-underground:permeability}{Doorlatend}: \omterm{def:g3:stones-pebbles-sand:sand}{zand}, \omterm{def:g3:stones-pebbles-sand:pebble}{grind}, krijt. \omterm{def:g6:water-underground:permeability}{Ondoorlatend}: klei, graniet zonder
scheuren, schalie.
\end{solution}

\begin{solution}{exo:g6:water-underground:2}
(a) \omterm{def:g6:water-underground:water-table}{grondwaterspiegel}; (b) \omterm{def:g6:water-underground:aquifer}{watervoerende laag}; (c) \omterm{def:g6:water-underground:spring}{bron}; (d) \omterm{def:g6:water-underground:spring}{put}.
\end{solution}

\begin{solution}{exo:g6:water-underground:3}
(a) Het verdampt weer in de lucht; het stroomt over het oppervlak weg naar
beken; het zakt in de grond (\omterm{def:g6:water-underground:infiltration}{infiltratie}). (b) Op de parkeerplaats: het
harde oppervlak laat bijna geen water in de grond zakken, dus stroomt het
water weg. In het weiland laten de bodem en het gras veel ervan in de grond
zakken.
\end{solution}

\begin{solution}{exo:g6:water-underground:4}
(a) Niet waar: klei is heel \omterm{def:g6:water-underground:porosity}{poreus} maar \omterm{def:g6:water-underground:permeability}{ondoorlatend}, omdat de \omterm{def:g6:water-underground:porosity}{poriën}
piepklein zijn. (b) Waar. (c) Niet waar: hij volgt de vorm van het land,
maar vlakker, en hij gaat met de seizoenen op en neer. (d) Waar: ongeveer
30~\% van al het zoete water is \omterm{def:g6:water-underground:groundwater}{grondwater}, tegen ongeveer een kwart
procent in de \omterm{def:g1:landscapes-around-us:water}{meren}.
\end{solution}

\begin{solution}{exo:g6:water-underground:5}
$10 \times 25 \div 100 = 2.5$: je kunt 2.5~liter erin gieten voordat hij
overloopt. $10 \times 22 \div 100 = 2.2$: er zou 2.2~liter door het gat
weglopen.
\end{solution}

\begin{solution}{exo:g6:water-underground:6}
(a) $2 \times 1 \times 1 = 2$~kubieke meter, dat is 2000~liter.
(b) $2000 \times 20 \div 100 = 400$~liter. (c) $2000 \times 18 \div 100 =
360$~liter loopt weg; $400 - 360 = 40$~liter blijft vastzitten.
\end{solution}

\begin{solution}{exo:g6:water-underground:7}
(a) De \omterm{def:g6:water-underground:spring}{bron} ligt op ongeveer 203~m, de \omterm{def:g1:landscapes-around-us:water}{rivier} op ongeveer 132~m. (b) De
\omterm{def:g6:water-underground:spring}{bron} ligt ongeveer 70~m boven de \omterm{def:g1:landscapes-around-us:water}{rivier}: het water stroomt de helling af
naar de \omterm{def:g1:landscapes-around-us:water}{rivier}. (c) Binnen in de krijthelling ligt de \omterm{def:g6:water-underground:water-table}{grondwaterspiegel}
onder de grond; hij raakt het oppervlak alleen waar de onderkant van het
krijt, die op de ondoorlatende schalie rust, op de helling naar buiten
komt.
\end{solution}

\begin{solution}{exo:g6:water-underground:8}
(a) Na het regenseizoen: $401 - 336 = 65$~m; na het droge seizoen:
$401 - 314 = 87$~m. (b) $336 - 314 = 22$~m.
\end{solution}

\begin{solution}{exo:g6:water-underground:9}
Dieper. De \omterm{def:g6:water-underground:water-table}{grondwaterspiegel} volgt de vorm van het land, maar vlakker:
boven op de \omterm{def:g1:landscapes-around-us:landform}{heuvel} ligt hij hoger dan onder \omterm{def:g6:water-underground:spring}{put}~A, maar de grond stijgt
veel meer, dus ligt het water verder onder het oppervlak (op de doorsnede
meer dan 110~m diep).
\end{solution}

\begin{solution}{exo:g6:water-underground:10}
De \omterm{def:g6:water-underground:porosity}{poriën} van klei zijn piepklein: het water blijft aan de schilfers
vastzitten en kan niet door de klei naar de \omterm{def:g6:water-underground:spring}{put} stromen. Klei is \omterm{def:g6:water-underground:porosity}{poreus}
maar \omterm{def:g6:water-underground:permeability}{ondoorlatend}, dus loopt er maar weinig water (20~liter van de 500)
uit weg.
\end{solution}

\begin{solution}{exo:g6:water-underground:11}
(a) $30.1 \div 0.26 \approx 116$: meer dan honderd keer zoveel.
(b) $100 - 69 = 31$: er blijft ongeveer 31~\% over, en \omterm{def:g6:water-underground:groundwater}{grondwater} alleen is
daar al ongeveer 30 van. (c) Het is veruit de grootste voorraad vloeibaar
\omterm{def:g2:water-on-earth:salt-fresh}{zoet water}; het ijs is bevroren en ligt meestal ver weg, in de poolgebieden
en op \omterm{def:g1:landscapes-around-us:landform}{bergen}. En \omterm{def:g6:water-underground:groundwater}{grondwater} zit bijna overal onder onze voeten.
\end{solution}

\begin{solution}{exo:g6:water-underground:12}
Beide putten halen water uit dezelfde \omterm{def:g6:water-underground:aquifer}{watervoerende laag} in het krijt. De
krachtige pomp haalde er sneller water uit dan de \omterm{def:g4:clouds-rain-snow:precipitation}{regen} het krijt bijvult,
dus zakte de \omterm{def:g6:water-underground:water-table}{grondwaterspiegel} in de hele omgeving, tot onder de bodem van
de tweede \omterm{def:g6:water-underground:spring}{put}, die droog kwam te staan. De boeren zouden kunnen afspreken
om minder te pompen, vooral in de zomer, samen één \omterm{def:g6:water-underground:spring}{put} te gebruiken, of
de tweede \omterm{def:g6:water-underground:spring}{put} dieper te maken; vooral moeten ze niet meer water opnemen
dan de \omterm{def:g4:clouds-rain-snow:precipitation}{regen} terugbrengt.
\end{solution}

\begin{solution}{pb:g6:water-underground:1}
\textbf{1.} De boerderij staat op krijt; daaronder ligt schalie. Het krijt
is \omterm{def:g6:water-underground:permeability}{doorlatend}, de schalie niet.

\textbf{2.} De \omterm{def:g4:clouds-rain-snow:precipitation}{regen} zakt in het krijt en stopt op de schalie. Waar de
onderkant van het krijt op de helling naar buiten komt, raakt de
\omterm{def:g6:water-underground:water-table}{grondwaterspiegel} de grond en komt het water naar buiten: dat is de \omterm{def:g6:water-underground:spring}{bron}.

\textbf{3.} $401 - 336 = 65$~m onder de grond.

\textbf{4.} De bodem van de \omterm{def:g6:water-underground:spring}{put} ligt op $401 - 100 = 301$~m, onder de
\omterm{def:g6:water-underground:water-table}{grondwaterspiegel} op 314~m: de \omterm{def:g6:water-underground:spring}{put} bereikt het water nog, met $314 - 301 =
13$~m.

\textbf{5.} $100 \times 50 \times 2 = 10\,000$~kubieke meter.

\textbf{6.} $10\,000 \times 20 \div 100 = 2000$~kubieke meter.

\textbf{7.} $10\,000 \times 18 \div 100 = 1800$~kubieke meter loopt weg;
$2000 - 1800 = 200$~kubieke meter blijft vastzitten.

\textbf{8.} $1800 \times 1000 = 1\,800\,000$~liter.

\textbf{9.} $18 \times 1000 = 18\,000$~liter per dag.

\textbf{10.} $18 \times 7 = 126$~kubieke meter per week.

\textbf{11.} De \omterm{def:g4:clouds-rain-snow:precipitation}{regen} van het natte seizoen zakt in het krijt en vult het
elk jaar weer aan (zolang de boerderijen niet meer oppompen dan de \omterm{def:g4:clouds-rain-snow:precipitation}{regen}
terugbrengt).

\textbf{12.} $1800 \div 18 = $~\textbf{100 dagen}: iets meer dan drie
maanden zonder enige \omterm{def:g4:clouds-rain-snow:precipitation}{regen}.
\end{solution}
